\chapter{Results and Discussion}
\label{chap:results}

This chapter presents the results in three parts that mirror the objectives
of Chapter~\ref{chap:introduction}: a physics validation and headline
benchmark (Section~\ref{sec:res_headline}), the data-efficiency study
(Section~\ref{sec:res_efficiency}), and the calibration study
(Section~\ref{sec:res_calibration}). The interpretability of the models is
examined in Section~\ref{sec:res_interpret} and the findings are summarised
in Section~\ref{sec:res_summary}.

\section{Physics Validation and Benchmark}
\label{sec:res_headline}

\subsection{The physics is correct}
Before comparing classifiers we confirm that the preprocessing reproduces
the expected physics. Figure~\ref{fig:substructure} shows the thirteen
engineered observables for signal and background. The top-jet mass peaks
near $m_t\simeq173$\,GeV while QCD jets are light; top jets have more
constituents and smaller $\tau_{32}$, the signature of three-prong structure.
Figure~\ref{fig:avg_images} shows the averaged jet images, in which the
characteristic three prongs of the top decay are visible in the
signal$-$background difference. These textbook features confirm that the
pipeline is physically sound.

\begin{figure}[htbp]
\centering
\includegraphics[width=\textwidth]{jet_substructure.pdf}
\caption{Distributions of the thirteen engineered jet-substructure
observables for top (signal) and QCD (background) jets on the test set. The
jet mass peaks at $m_t$ for top jets; $\tau_{32}$ and the constituent
multiplicity also separate the classes, as expected from physics.}
\label{fig:substructure}
\end{figure}

\begin{figure}[htbp]
\centering
\includegraphics[width=\textwidth]{average_jet_images.pdf}
\caption{Average $p_T$-weighted jet images for top jets (left), QCD jets
(centre) and their difference (right). The three-prong structure of the
hadronic top decay is visible in the difference panel.}
\label{fig:avg_images}
\end{figure}

\subsection{Model comparison}
Table~\ref{tab:model_comparison} reports the headline performance of the four
models, and Figure~\ref{fig:roc} shows the ROC curves and background
rejection versus signal efficiency. The Particle Transformer achieves the
best discrimination (AUC $=0.978$, background rejection
$1/\varepsilon_B\approx136$ at $\varepsilon_S=0.5$), followed by ParticleNet,
the CNN and the BDT. This reproduces the ordering reported in the community
comparison~\citep{kasieczka2019landscape}; the absolute values sit slightly
below the published saturation values because we train on a subset rather
than the full $1.2\times10^{6}$ jets.

\IfFileExists{../../../results/comparison_table.tex}
  {\input{../../../results/comparison_table.tex}}
  {\emph{(comparison table generated by the pipeline)}}

\begin{figure}[htbp]
\centering
\includegraphics[width=\textwidth]{roc_comparison.pdf}
\caption{Left: ROC curves for the four models on the test set. Right:
background rejection $1/\varepsilon_B$ as a function of signal efficiency
$\varepsilon_S$. The Particle Transformer (red) is best across the range.}
\label{fig:roc}
\end{figure}

Figure~\ref{fig:scores} shows the classifier-score distributions and
Figure~\ref{fig:confusion} the confusion matrices at a threshold of $0.5$;
both confirm clean signal/background separation for all models, sharpest for
the deep architectures.

\begin{figure}[htbp]
\centering
\includegraphics[width=\textwidth]{score_distributions.pdf}
\caption{Classifier-score distributions for top (red) and QCD (blue) jets.
Well-separated peaks near $0$ and $1$ indicate strong discrimination.}
\label{fig:scores}
\end{figure}

\begin{figure}[htbp]
\centering
\includegraphics[width=\textwidth]{confusion_matrices.pdf}
\caption{Confusion matrices (normalised) at a decision threshold of $0.5$.}
\label{fig:confusion}
\end{figure}

\section{Data-Efficiency Study}
\label{sec:res_efficiency}

Figure~\ref{fig:learning} shows the test AUC and background rejection as a
function of the training-set size $N_{\text{train}}$. Three findings stand
out:
\begin{enumerate}
    \item \textbf{The BDT is the most data-efficient model.} Its AUC rises by
    under one percent (from $0.961$ to $0.968$) across the whole
    $5$k--$50$k range; engineered features short-circuit much of what the
    deep models must learn from data.
    \item \textbf{The deep models show a clear crossover.} The Particle
    Transformer is the \emph{weakest} model at $5$k jets (AUC $=0.949$) but
    the \emph{strongest} at $50$k ($0.974$), overtaking both the BDT and
    ParticleNet. Its slope is still positive at the largest size tested,
    consistent with the literature observation that it only saturates near
    $10^{6}$ jets~\citep{qu2022particletransformer}.
    \item \textbf{The CNN is intermediate}, gaining substantially with data
    ($0.947\rightarrow0.968$) but plateauing earlier than the
    permutation-equivariant graph and attention models.
\end{enumerate}
The practical message is that under a tight Monte~Carlo budget the simple,
interpretable BDT is competitive or better, whereas the attention model only
realises its advantage once enough data is available.

\begin{figure}[htbp]
\centering
\includegraphics[width=\textwidth]{learning_curves.pdf}
\caption{Data-efficiency curves. Left: test AUC versus training-set size with
$1\sigma$ bootstrap error bars. Right: background rejection at
$\varepsilon_S=0.5$. The Particle Transformer rises from worst at $5$k jets
to best at $5\times10^{4}$, crossing the other models.}
\label{fig:learning}
\end{figure}

\section{Probability-Calibration Study}
\label{sec:res_calibration}

Figure~\ref{fig:calibration} shows the reliability diagrams and score
histograms, and Table~\ref{tab:calibration} reports the ECE, the fitted
temperature $T$ and the post-rescaling ECE. The key results are:
\begin{enumerate}
    \item \textbf{All models are mildly over-confident} (every fitted
    temperature satisfies $T>1$), in line with
    \citet{guo2017calibration}.
    \item \textbf{Ranking and calibration are decoupled.} The jet-image CNN
    is the worst-calibrated model by an order of magnitude (ECE $=0.078$,
    $T=1.54$) despite ranking only third on AUC, whereas the top-ranked
    Particle Transformer is comparatively well-calibrated (ECE $=0.012$). The
    BDT and ParticleNet are already well-calibrated (ECE $\approx0.007$).
    \item \textbf{Temperature scaling helps, but not uniformly.} A single
    scalar brings the BDT, ParticleNet and Particle Transformer to
    ECE\,$\lesssim0.008$, but only halves the CNN's error (to $0.064$): the
    CNN's miscalibration is structural and would need a richer recalibration
    such as isotonic regression.
\end{enumerate}

\begin{figure}[htbp]
\centering
\includegraphics[width=\textwidth]{calibration.pdf}
\caption{Top: reliability diagrams for the four models (solid: raw scores;
dotted: after temperature scaling). Bottom: score histograms. The jet-image
CNN deviates most strongly from the diagonal.}
\label{fig:calibration}
\end{figure}

\begin{table}[htbp]
\centering
\caption{Test-set calibration at the $5\times10^{4}$-jet study endpoint. $T$
is the temperature fitted on the validation split; $\mathrm{ECE}_T$ is the
Expected Calibration Error after rescaling.}
\label{tab:calibration}
\ra{1.2}
\IfFileExists{../../../results/calibration_table.tex}
  {\input{../../../results/calibration_table.tex}}
  {\emph{(calibration table generated by the pipeline)}}
\end{table}

\section{Model Interpretability}
\label{sec:res_interpret}

Two views of \emph{why} the models work are useful. Figure~\ref{fig:bdt_imp}
shows the BDT feature importances: the jet mass (gain $118$), the
$1$-subjettiness $\tau_1$ ($89$) and the jet width ($76$) dominate, followed
by the constituent multiplicity ($15$) --- all sensitive to the larger mass,
girth and prong count of top jets. The canonical ratio $\tau_{32}$ ranks
lower (gain $7$): the boosted trees capture the prong information mainly
through the absolute $\tau_1$ and the jet width rather than the ratio, even
though $\tau_{32}$ is the cleaner discriminant at the distribution level
(Fig.~\ref{fig:substructure}).
Figure~\ref{fig:cls_attn} shows the class-token attention of the Particle
Transformer versus a constituent's angular distance $\Delta R$ from the jet
axis. For QCD jets the attention concentrates at small $\Delta R$ (a single
core), whereas for top jets it is spread to much larger $\Delta R$, reflecting
the broader, multi-pronged energy flow of the top decay. The model has
learned to attend to the peripheral radiation that carries the top signature.

\begin{figure}[htbp]
\centering
\includegraphics[width=0.8\textwidth]{bdt_importance.pdf}
\caption{Gain-based feature importance of the BDT. The jet mass, $\tau_1$
and jet width dominate, followed by the constituent multiplicity.}
\label{fig:bdt_imp}
\end{figure}

\begin{figure}[htbp]
\centering
\includegraphics[width=\textwidth]{cls_attention.pdf}
\caption{Class-token attention weight of the Particle Transformer versus the
constituent angular distance $\Delta R$ from the jet axis, for top jets
(left) and QCD jets (right). Top jets draw attention to larger $\Delta R$,
matching their multi-prong structure.}
\label{fig:cls_attn}
\end{figure}

\section{Summary of Findings}
\label{sec:res_summary}

\begin{enumerate}
    \item The pipeline reproduces the established ordering Particle
    Transformer $>$ ParticleNet $>$ CNN $>$ BDT, with the Particle
    Transformer reaching AUC $=0.978$.
    \item Data efficiency separates the models sharply: the BDT is best in
    the small-data regime, while the deep models overtake it as data grows,
    producing a clear crossover.
    \item Ranking power and calibration are decoupled: the CNN is the
    worst-calibrated model despite a mid-pack ranking, and temperature
    scaling fixes every model except the CNN.
\end{enumerate}
